\section{CDF Mapping (Experimental)}
\label{sec:cdf}

An experimental approach: use the empirical CDF of the data as the locality-preserving
hash $h$, redistributing the FPR budget from sparse regions to dense regions.

\textbf{Status:} theoretical exploration. Benchmarks did not show practical improvement
over the SODA hash or hybrid segmentation.

\subsection{Motivation}

The information-theoretic lower bound (Theorem~1) holds for all data and query
distributions. But if queries follow approximately the same distribution as stored
keys --- users query ranges where data exists, not empty regions --- then FPR in
sparse inter-cluster gaps does not matter. If we could \emph{move} FPR budget away from
dense regions and into sparse regions, the effective FPR experienced by real queries
would be lower for the same number of bits.

\subsection{The Hash Function}

The CDF is a monotonic function that maps keys to a near-uniform distribution
(Figure~\ref{fig:cdf-redistribution}):
\begin{itemize}
  \item \textbf{Where keys are dense} (clusters): the mapped space is stretched, so
    the same number of non-keys is spread over a wider interval $\Rightarrow$ fewer
    non-keys per unit length $\Rightarrow$ lower local FPR.
  \item \textbf{Where keys are sparse} (gaps): the mapped space is compressed, non-keys
    pile up $\Rightarrow$ higher local FPR. But under our assumption, queries rarely
    land here.
\end{itemize}

Monotonicity ($x_1 < x_2 \Rightarrow h(x_1) \leq h(x_2)$) preserves range queries ---
intervals map to intervals.

\begin{figure}[h]
  \centering
  \includegraphics[width=0.9\textwidth]{cdf_redistribution}
  \caption{CDF redistribution. Top: original universe with non-uniform key density
    (uniform background + two clusters). After CDF mapping: keys become near-uniform,
    non-key density inverts --- valleys where clusters were (low FPR), peaks where
    gaps were (high FPR, but queries rarely land there).}
  \label{fig:cdf-redistribution}
\end{figure}

\begin{figure}[h]
  \centering
  \includegraphics[width=0.9\textwidth]{cdf_mapping}
  \caption{Before and after CDF mapping. Dense clusters (top) are stretched into
    uniform blocks (bottom); sparse gaps are compressed.}
  \label{fig:cdf-mapping}
\end{figure}

\textbf{Important caveat:} because the CDF maps everything monotonically, keys and
non-keys are still interleaved in $U'$ --- there is no strict separation. The FPR
reduction is not formally bounded and depends entirely on the query distribution matching
the data distribution.

\subsection{ARE Construction With CDF Hash}

\begin{enumerate}
  \item Sort input keys.
  \item Build a PGM index~\citep{ferragina2020pgm} on float64-converted keys to obtain
    an approximate CDF.
  \item Fix monotonicity (PGM positions may be non-monotone due to float64 rounding)
    via running max.
  \item Sample CDF control points every $\texttt{pgmEpsilon}$ keys $\Rightarrow$
    piecewise-linear model.
  \item Compute $L_\text{eff}$: the effective range length after CDF distortion.
    A query of length $\mathcal{L}$ maps to a range of variable width depending on
    the local CDF slope, bounded by
    $L_\text{eff} = \mathcal{L} \times \max_\text{segment} \text{density}$.
  \item Set $K = \lceil \log_2(n \cdot L_\text{eff} / \varepsilon) \rceil$.
  \item Map keys through the piecewise-linear CDF, deduplicate, build ERE.
  \item Query: map both endpoints through the same CDF, forward to ERE.
\end{enumerate}

CDF model cost: each control point stores a (key, rank) pair $= 128$ bits. With
$\texttt{pgmEpsilon} = 128$: ${\approx} n/128$ points $\approx 1$ bit per key.

\subsection{Limitations}

\begin{itemize}
  \item $O(n^2)$ build time due to PGM convex hull construction; constructor returns
    error for $n > 2^{20}$.
  \item Float64 precision loss for keys $> 2^{53}$.
  \item FPR is not formally bounded --- depends on query/data distribution match.
  \item CDF model overhead adds 1--2 BPK on top of ERE storage.
\end{itemize}

\subsection{Experimental Results}

Benchmarks against the SODA hash showed that CDF-ARE does not provide a consistent
practical advantage:
\begin{itemize}
  \item When query distribution matches data: CDF-ARE uses fewer bits per key, but FPR
    is comparable or worse due to $L_\text{eff}$ amplification in dense regions.
  \item When query distribution diverges from data: FPR degrades catastrophically
    (40--58\% at $\sigma_\text{query} = 4 \times \sigma_\text{data}$).
  \item The $O(n^2)$ build cost makes it impractical for large datasets.
\end{itemize}

Hybrid segmentation (Chapter~\ref{chap:hybrid}) achieves a similar goal --- exploiting
clustering --- more effectively via segmentation and exact mode, without the CDF overhead.
